\section{Endogenous pension design}\label{app:esc}\setcounter{equation}{0}
This appendix collects the derivations behind section \ref{sec:esc}: the corner that costless redistribution produces, the deadweight cost of redistribution, what the cost changes in the equilibrium of appendices \ref{app:EE} and \ref{app:PEE}, and the state independence of the design choice under log preferences. Throughout there are no informal households, $\gamma_0 = 0$, so the budget \eqref{eq:model:governmentBudget} gives $\overline{b}_t = \nu_t w_t h_t\tau_t/h_{t-1}$ and a retiree of type $i$ the benefit $\nu_t w_t h_t\tau_t\left(1+\theta_t\left[y_i-1\right]\right)$, where $y_i \equiv \eta_i^{1+\xi}/(X_i^{\xi}\Gamma_h) = \eta_i h_{t-1,i}/h_{t-1}$ is the labor income of type $i$ relative to the mean by \eqref{eq:hi_h}, and $\sum_i\gamma_i y_i = 1$.

\subsection{Costless redistribution}\label{app:esc:corner}
We consider three timings of the choice. If $\theta_t$ is chosen at $t$ together with $\tau_t$, it re-splits benefits that current retirees have already earned. Current labor supply and savings depend on the design that will apply to the benefits they earn, $\theta_{t+1}$, so the choice moves no economic decision, and its marginal effect on the political objective \eqref{eq:LOG:PolObj} runs entirely through current retirees,
\begin{align}\label{eq:esc:seqFOC}
   \dfrac{\partial \mathcal{W}_t}{\partial \theta_t} =\omega \sum_i\gamma_i\mu_i\frac{\frac{1-\alpha}{\alpha}\tau_t\left(y_i-1\right)}{\frac{s_{t-1}^i}{s_{t-1}}+\frac{1-\alpha}{\alpha}\tau_t\left(1+\theta_t\left[y_i-1\right]\right)}.
\end{align}
The numerator is the change in the benefit of type $i$ and averages to zero across types: a more Bismarckian split moves a given total from retirees below mean income to retirees above it. The denominator is proportional to the consumption of a retiree of type $i$ and rises with income. With equal propensities to vote the sum therefore weights the gains of the rich by a lower marginal utility than the losses of the poor. The following proposition makes the argument exact.

\begin{proposition}\label{prop:esc:corner}
Assume log-GHH preferences, no informal households and no deadweight cost of redistribution (section \ref{sec:esc:cost}), and let the design $\theta_t$ that applies to the benefits paid at $t$ be chosen at $t$ together with $\tau_t$. Let the propensities to vote be positive and not rise with relative income, $\mu_i \leq \mu_j$ whenever $y_i > y_j$, equal propensities included, and let relative incomes not all be equal. Then, for every $\omega > 0$ and $\tau_t > 0$, the derivative \eqref{eq:esc:seqFOC} is negative at every $\theta_t \in [0,1]$ at which every retiree's consumption is positive, and the design chosen together with a positive tax is the Beveridgean corner, $\theta_t = 0$. This holds for all values of $\alpha$, $\beta$, $\xi$, $\nu_t$, $\gamma_i$, $\eta_i$ and $X_i$, and whatever tax and design current retirees expected when they saved.
\end{proposition}

\smalltitle{Proof of proposition \ref{prop:esc:corner}.} By \eqref{eq:LOG:formalSol_i} the labor supply and the savings of the young at $t$ depend on the design that will apply to their own benefits, $\theta_{t+1}$, and not on $\theta_t$, so $\theta_t$ enters no state at $t+1$ and no later decision. It enters $\mathcal{W}_t$ only through the consumption of current retirees (appendix \ref{app:EE:CRRA}, with $\gamma_0 = 0$),
\begin{align}\label{eq:esc:corner:c2}
    c_{2,t}^i = \alpha\nu_t h_t^{1-\alpha}\left(\frac{s_{t-1}}{\nu_t}\right)^{\alpha} D_i, \qquad D_i \equiv \frac{s_{t-1}^i}{s_{t-1}}+\frac{1-\alpha}{\alpha}\tau_t\left(1+\theta_t\left[y_i-1\right]\right),
\end{align}
where the first factor is common to all types and independent of $\theta_t$. From \eqref{eq:LOG:PolObj},
\begin{align*}
    \frac{\partial \mathcal{W}_t}{\partial \theta_t} = \omega\sum_i\gamma_i\mu_i\frac{\partial \ln c_{2,t}^i}{\partial \theta_t} = \omega\,\frac{1-\alpha}{\alpha}\,\tau_t\sum_i\gamma_i\left(y_i-1\right)\frac{\mu_i}{D_i},
\end{align*}
which is \eqref{eq:esc:seqFOC}. We sign the sum in three steps.

First, relative incomes average to one. By the definition of $\Gamma_h$ in \eqref{eq:hi_h}, $\sum_i\gamma_i y_i = 1$, and since $\sum_i\gamma_i = 1$ (section \ref{sec:model:Demo}), $\sum_i\gamma_i\left(y_i-1\right) = 0$: the design moves a given total between retirees below and above mean income.

Second, $D_i$ rises with $y_i$. The savings of current retirees were chosen at $t-1$. By part iii of proposition \ref{prop:LOG:EE}, that is \eqref{eq:si_s} with $B_{t+1} = \beta$, in which the tax and the design are those that apply to the benefits of the generation that saves,
\begin{align*}
    \frac{s_{t-1}^i}{s_{t-1}} = y_i + \frac{1-\alpha}{\alpha}\,\frac{\tau^{e}_t\left(1-\theta^{e}_t\right)}{1+\beta}\left(y_i-1\right),
\end{align*}
where $\tau^{e}_t$ and $\theta^{e}_t$ are the tax and the design that generation $t-1$ expected for period $t$. Whatever these were, the slope in $y_i$ is at least one, since $\tau^{e}_t \geq 0$ and $\theta^{e}_t \leq 1$. The second term of $D_i$ has slope $\frac{1-\alpha}{\alpha}\tau_t\theta_t \geq 0$, so $D_i$ rises with $y_i$ with a slope of at least one, and by \eqref{eq:esc:corner:c2} it is positive wherever every retiree's consumption is.

Third, Chebyshev's sum inequality: for weights $\gamma_i \geq 0$ that sum to one,
\begin{align*}
    \sum_i\gamma_i a_i b_i - \sum_i\gamma_i a_i\sum_j\gamma_j b_j = \frac{1}{2}\sum_{i,j}\gamma_i\gamma_j\left(a_i-a_j\right)\left(b_i-b_j\right),
\end{align*}
which is non-positive when $a_i$ and $b_i$ are ordered in opposite directions. Let $a_i = y_i-1$ and $b_i = \mu_i/D_i$. If $y_i > y_j$, then $\mu_i \leq \mu_j$ and $D_i > D_j > 0$, so $b_i < b_j$, and every term with $y_i \neq y_j$ is negative. Since $\sum_i\gamma_i a_i = 0$ by the first step and relative incomes are not all equal, $\sum_i\gamma_i\left(y_i-1\right)\mu_i/D_i < 0$, and so is the derivative.

Finally, $D_i$ falls with $\theta_t$ only for types below mean income, and $D_i \geq 1$ for the others. The designs at which every retiree's consumption is positive therefore form an interval that starts at $\theta_t = 0$ whenever there are any, $\mathcal{W}_t$ is strictly decreasing on it, and the design chosen with any positive tax is $\theta_t = 0$.

The proof uses the propensities to vote only through the ratios $\mu_i/D_i$, which order the types as $\mu_i/c_{2,t}^i$ does. The derivative is therefore negative at every design at which propensities to vote rise with relative income, if at all, by proportionally less than consumption in retirement, $\mu_i/\mu_j < c_{2,t}^i/c_{2,t}^j$ whenever $y_i > y_j$, and a design other than the Beveridgean corner requires propensities to vote that, for some pair of types and at some design, rise at least as fast as consumption. The argument also holds under CRRA-GHH preferences at every intertemporal elasticity: the design still enters only the consumption of current retirees, whose marginal utility is $D_i^{-1/\rho}$ times a factor common to all types, and relative savings keep the form of \eqref{eq:si_s} with a discount term $B_t$ common to all types, so $\mu_i D_i^{-1/\rho}$ takes the place of $\mu_i/D_i$.

Propensities to vote that rise with income can overturn the corner, and they do rise in the U.S.\ calibration of section \ref{sec:oecd}, so there the sign is a quantitative matter: evaluated at the calibrated 2020 equilibrium, \eqref{eq:esc:seqFOC} is negative over the whole interval, at every date of the projection and at $\rho = 0.5$, $1$ and $2$, and the choice is the corner $\theta_t = 0$.

If instead $\theta_{t+1}$ is chosen at $t$, the design no longer touches the benefits of current retirees but moves current labor supply and savings, directly and through the anticipated response of $\tau_{t+1}$. With the political objective \eqref{eq:LOG:PolObj} the marginal effect is
\begin{align}\label{eq:esc:leadFOC}
   \dfrac{\partial \mathcal{W}_t}{\partial \theta_{t+1}} = \sum_i\gamma_i\mu_i\left[\omega\,(1-\alpha)\dfrac{d\ln\Theta_{h,t}}{d\theta_{t+1}} + \nu_t\left((1+\beta)\dfrac{d\ln\tilde{c}_{1,t}^i}{d\theta_{t+1}} + \beta\dfrac{d\ln R_{t+1}}{d\theta_{t+1}}\right)\right],
\end{align}
where $\Theta_{h,t}$ is the labor supply function \eqref{eq:h} and each total derivative includes the response of $\tau_{t+1}$, and of $\theta_{t+2}$, to the design chosen. Retirees are unanimous, since their term does not depend on $i$, and they favour a more contributive system: a higher $\theta_{t+1}$ raises the return to an hour worked at $t$ and with it the labor supply on which their benefits are levied. The young are divided between this forward-looking stake and the redistribution of their own future benefits, the force of \eqref{eq:esc:seqFOC} one period ahead. The equilibrium is a sequence of policy functions over the inherited design, identified by backward induction as in section \ref{sec:numerical}, and at observed U.S.\ inequality the redistributive force dominates: the choice is again the corner $\theta = 0$.

A design chosen once and for all enters every channel at once. The joint choice of $\tau_{t_0}$ and $\theta$ is nevertheless one-dimensional, since at any candidate design the first order condition for the tax is the tax condition of appendix \ref{app:PEE:LOG} at that design, and the predetermined distribution of savings is held at the value the chosen design implies, which makes the equilibrium a fixed point. The outcome is again a corner: $\theta = 0$ under log preferences, and under CRRA preferences the Bismarckian corner $\theta = 1$ at high intertemporal elasticities, where the forward-looking channels are strong. Near the elasticity at which the corner switches the electorate is close to indifferent between the two, so the switch is not sharply located. The online appendix reports the three timings at each intertemporal elasticity (\oa{esc-timing}); the technical documentation gives the algorithms \parencite{BergGE26doc}.

\subsection{The deadweight cost of redistribution}\label{app:esc:cost}
\smalltitle{The implicit tax.} A retiree of type $i$ receives $\nu_t w_t h_t\tau_t\left[\theta_t y_i + 1-\theta_t\right]$ on past earnings $w_{t-1}\eta_i h_{t-1,i} = w_{t-1}y_i h_{t-1}$, that is
\begin{align*}
    \Pi_t\,\tau_t\left[\theta_t + \frac{1-\theta_t}{y_i}\right] \text{ per unit of past earnings}, \qquad \Pi_t \equiv \frac{\nu_t w_t h_t}{w_{t-1}h_{t-1}},
\end{align*}
where $\Pi_t$ is the pay-as-you-go return factor: a purely earnings-related system paying the same revenue returns $\Pi_t\tau_t$ per unit of earnings to every type. The flat component moves type $i$ away from that benchmark by
\begin{align}\label{eq:esc:implicitTax}
    t_i \equiv \Pi_t\,\tau_t\,(1-\theta_t)\left(1-\frac{1}{y_i}\right)
\end{align}
per unit of earnings: a tax on types above mean income, a subsidy to types below it, and a pure transfer in the aggregate, since $\sum_i\gamma_i y_i t_i = 0$. This is the tax component of the contribution in the sense of \textcite{Summers89a1}, the part of the payroll charge that buys the contributor no benefit, and the only part that distorts like a tax.

\smalltitle{Which margin.} The wedge \eqref{eq:esc:implicitTax} is an average implicit tax: it compares what type $i$ receives per unit of its earnings with the benchmark. The marginal wedge is a different object. The flat component depends on aggregate hours and not on own hours, so an extra hour raises the benefit of every type by $\Pi_t\tau_t\theta_t$ per unit of earnings, and the marginal wedge of the flat component, $\Pi_t\tau_t(1-\theta_t)$, is common to all types. That wedge is already in the model: it is the group $\frac{1-\alpha}{\alpha}\theta\tau_{t+1}\Gamma_{s,t}$ in the labor supply function \eqref{eq:h}, through which hours respond to the design with the Frisch elasticity $\xi$, and GHH preferences carry no income effect through which the transfers between types could move hours. A cost built on the hours margin would count that response twice and could not depend on the spread of incomes. The cost therefore stands for margins the model does not contain, participation, formality and the reporting of earnings, on which the average wedge is the relevant one: a worker who decides whether to bring earnings into the system compares the earnings at stake with what they buy, so the wedge that matters is the average tax on those earnings \parencite{Saez02a1,KlevenK06a1}. The evidence finds the tax component to matter on these margins. \textcite{Disney04a1} splits pension contributions in OECD countries into a tax and a saving component and finds that the tax component lowers the participation of women while the saving component raises it, and \textcite{KumlerVF20a1} find that the underreporting of wages to Mexico's social security agency fell for the workers whose benefits the 1997 reform tied more closely to reported wages.

\smalltitle{The Harberger loss.} A set of proportional wedges $t_i$ on bases $e_i$, each responding to its net-of-tax rate with a common elasticity $\varepsilon$, costs $\frac{1}{2}\varepsilon\sum_i\gamma_i t_i^2 e_i$ to second order in the wedges \parencite{Harberger64a1}. Per unit of mean past earnings the bases are $e_i = y_i$, so the loss is $\frac{1}{2}\varepsilon\,\Pi_t^2\tau_t^2(1-\theta_t)^2\tilde{V}$ with
\begin{align}\label{eq:esc:Vtilde}
    \tilde{V} \equiv \sum_i\gamma_i\frac{\left(y_i-1\right)^2}{y_i},
\end{align}
while revenue per unit of mean past earnings is $\Pi_t\tau_t$. As a share of revenue the loss is therefore $\frac{1}{2}\varepsilon\,\Pi_t\,\tau_t\,(1-\theta_t)^2\,\tilde{V}$: quadratic in the flat share of the system, linear in its size and proportional to a dispersion of relative incomes that is zero when all incomes are equal. The quadratic counts the subsidy to types below mean income as a loss, as it does the tax on types above it; where participation is already taxed, lowering the participation tax of low-income workers can reduce the total loss \parencite{Saez02a1}, so the symmetry is a property of the reduced form rather than of the argument.

\smalltitle{Specification.} The exponential form $f(\theta_t,\tau_t)$ of \eqref{eq:esc:budget} scales both components of the benefit alike,
\begin{align}\label{eq:esc:AB}
    A(\theta_t,\tau_t) = f(\theta_t,\tau_t)\,\theta_t, \qquad B(\theta_t,\tau_t) = f(\theta_t,\tau_t)\,(1-\theta_t),
\end{align}
so that the surviving share of revenue is $A+B = f \in (0,1]$ and $1-f$ is, to first order, the Harberger share with $\lambda$ in the place of $\varepsilon\Pi_t$.\footnote{The return factor multiplies every implicit tax alike and so scales the loss once. It carries $\nu_t$ and would make a structural cost parameter drift with demographics; we absorb it into $\lambda$, which is calibrated rather than derived.} The exponential keeps $f$ positive for every $\lambda$ and every size of system, and nests the model without the cost at $\lambda = 0$. Three properties follow. Nothing is lost when nothing is redistributed, $f = 1$ at $\theta_t = 1$ and for every design when $\tilde{V} = 0$. The loss per unit of redistribution rises with the size of the system, through $\tau_t$. And the marginal cost of the first unit of redistribution is zero, $\partial_{\theta}\ln f = \lambda\tau_t\tilde{V}(1-\theta_t)$ vanishes at $\theta_t = 1$, so under a purely quadratic loss no electorate chooses the Bismarckian corner, which is why France's design has no own cost in section \ref{sec:esc:calibration}.

\smalltitle{Magnitude.} The shape of $f$ follows from the tax component; its size does not. With $\Pi_t$ of the order of $\nu_t$, $1.34$ in the U.S.\ calibration for 2020 (table \ref{table:US:Calib}), participation and taxable-income elasticities of $0.1$ to $0.4$ \parencite{ChettyGMW11a1,SaezSG12a1} correspond to $\lambda$ between about $0.13$ and $0.54$. The calibrated values of table \ref{table:US_ESC:calibration} are more than an order of magnitude above that range at $\rho \leq 1$, and three to thirteen times above it at $\rho = 2$, so $\lambda$ also stands in for administrative, evasion and political costs of redistribution that the model does not contain.

\subsection{The equilibrium with the cost}\label{app:esc:equilibrium}
With the cost, the benefit formula of \eqref{eq:model:governmentBudget} becomes
\begin{align}\label{eq:esc:benefits}
    b_t^i = \left[A(\theta_t,\tau_t)\,h_{t-1,i}\eta_i + B(\theta_t,\tau_t)\,h_{t-1}\right]\overline{b}_t,
\end{align}
with $\overline{b}_t = \nu_t w_t h_t\tau_t/h_{t-1}$ the gross revenue per unit of past labor, of which a share $1-f$ does not reach retirees. The model without the cost is $(A,B) = (\theta_t, 1-\theta_t)$.

\smalltitle{Economic equilibrium.} In appendix \ref{app:EE} the design enters only through the groups $\frac{1-\alpha}{\alpha}\tau_{t+1}\theta_{t+1}$ and $\frac{1-\alpha}{\alpha}\tau_{t+1}(1-\theta_{t+1})$, in $\Gamma_{s,t}$, $\Theta_{h,t}$ and $s_{t,i}/s_t$, and through $\theta_t$ and $1-\theta_t$ in the benefits of current retirees. The cost is therefore exactly the substitution $\theta_{t+1} \mapsto A(\theta_{t+1},\tau_{t+1})$ and $1-\theta_{t+1} \mapsto B(\theta_{t+1},\tau_{t+1})$ in those groups, and $\theta_t \mapsto A(\theta_t,\tau_t)$, $1-\theta_t \mapsto B(\theta_t,\tau_t)$ in the benefit formula, and it changes nothing else. Under log preferences, \eqref{eq:s}, \eqref{eq:h} and \eqref{eq:si_s} read
\begin{subequations}\label{eq:esc:auxiliary}
\begin{align}
    \Gamma_{s,t} &= \frac{1}{1+\xi}\,\frac{\Gamma_h\beta}{1+\beta+\frac{1-\alpha}{\alpha}\tau_{t+1}\left[A(\theta_{t+1},\tau_{t+1})(1+\beta) + B(\theta_{t+1},\tau_{t+1})\right]}, \label{eq:esc:auxiliary:Gammas}\\
    \Theta_{h,t} &= \left(\Gamma_{h}\right)^{\frac{1+\xi}{1+\alpha\xi}}\left(\frac{(1-\alpha)(1-\tau_t)}{\Gamma_{h}-\frac{1-\alpha}{\alpha}A(\theta_{t+1},\tau_{t+1})\,\tau_{t+1}\Gamma_{s,t}}\right)^{\frac{\xi}{1+\alpha\xi}}, \label{eq:esc:auxiliary:Thetah}\\
    \frac{s_{t,i}}{s_t} &= y_i + \frac{1-\alpha}{\alpha}\,\frac{B(\theta_{t+1},\tau_{t+1})\,\tau_{t+1}}{1+\beta}\left(y_i-1\right), \label{eq:esc:auxiliary:si}
\end{align}
\end{subequations}
and the consumption of a retiree of type $i$ carries $A(\theta_t,\tau_t)\,y_i + B(\theta_t,\tau_t)$ in place of $1+\theta_t\left[y_i-1\right]$ in \eqref{eq:esc:corner:c2}. The two coincide when $A+B = 1$, so the model without the cost is nested exactly. Here $\tau_{t+1}$ is the tax the young anticipate for the period in which their benefits are paid, $\tau_{t+1} = \mathcal{T}_{t+1}(\theta_{t+1})$ with $\mathcal{T}_{t+1}$ the equilibrium tax function below, so the young anticipate both the design and the size of the system that will pay them.

\smalltitle{The tax condition.} The first order condition for the tax of appendix \ref{app:PEE:LOG} picks up one factor. The benefit term of the retirees' marginal effect $\mathcal{E}_{2,t}^i$ is multiplied by $\tau_t$, and $f$ now moves with $\tau_t$, so
\begin{align}\label{eq:esc:E2}
    \mathcal{E}_{2,t}^i = (1-\alpha)\frac{\partial \ln\left(\Theta_{h,t}\right)}{\partial \tau_t}+\frac{\frac{1-\alpha}{\alpha}\left[A(\theta_t,\tau_t)\,y_i + B(\theta_t,\tau_t)\right]\left(1+\tau_t\,\partial_{\tau}\ln f(\theta_t,\tau_t)\right)}{\frac{s_{t-1}^i}{s_{t-1}}+\frac{1-\alpha}{\alpha}\tau_t\left[A(\theta_t,\tau_t)\,y_i + B(\theta_t,\tau_t)\right]},
\end{align}
where $\partial_{\tau}\ln f = -\tfrac{1}{2}\lambda\tilde{V}(1-\theta_t)^2$ and hence $1+\tau_t\,\partial_{\tau}\ln f = 1+\ln f(\theta_t,\tau_t)$, since $\ln f$ is linear in $\tau_t$: a larger system is worth less to a retiree at the margin than the revenue it collects, by the share of the marginal contribution that the flat component dissipates. Nothing else in the condition sees $f$. The term of the young, $\mathcal{E}_{1,t}^i$, contains no policy but $\tau_t$, and $\partial\ln(\Theta_{h,t})/\partial\tau_t = -\frac{\xi}{1+\alpha\xi}\frac{1}{1-\tau_t}$ is free of the groups, so the condition involves $\theta_{t+1}$ nowhere and the equilibrium tax at $t$ is a function $\mathcal{T}_t(\theta_t)$ of the design in force and of parameters and demographics alone, as in proposition \ref{prop:LOG:PEE}.

\smalltitle{Identification of $\theta$.} The replacement rate of type $j$, its benefit over its past earnings, is $\Pi_t\tau_t\left[A + B/y_j\right]$. The ratio of the replacement rates of two types, the datum that identifies $\theta$ in section \ref{sec:oecd}, therefore depends on $(A,B)$ only through $B/A = (1-\theta_t)/\theta_t$, from which $f$ cancels: the design is identified separately from the cost, and $\theta^{\ast} = 0.738$ for the U.S.\ whatever $\lambda$. This is a property of the proportional specification. A cost on the flat component alone, $(A,B) = (\theta_t, f(1-\theta_t))$, leaves $B/A = f(1-\theta_t)/\theta_t$, so that the same datum implies a different $\theta$ for every value of the cost parameter and the two are identified jointly.

\smalltitle{The design condition.} With the design chosen one period in advance, $\theta_{t+1}$ enters the objective through $A(\theta_{t+1},\tau_{t+1})$ and $B(\theta_{t+1},\tau_{t+1})$ in \eqref{eq:esc:auxiliary}, with $\tau_{t+1}$ the anticipated equilibrium tax. At a given tax, $\partial_{\theta}A = f + \theta_{t+1}f_{\theta}$ and $\partial_{\theta}B = (1-\theta_{t+1})f_{\theta} - f$ with $f_{\theta} = f\lambda\tau_{t+1}\tilde{V}(1-\theta_{t+1})$, so
\begin{align}\label{eq:esc:ftheta}
    \partial_{\theta}A + \partial_{\theta}B = f_{\theta} = f\,\lambda\,\tau_{t+1}\,\tilde{V}\,(1-\theta_{t+1}) \geq 0,
\end{align}
zero at $\theta_{t+1} = 1$ and rising with the flat share, the dispersion of incomes and the size of the system the young anticipate. Without the cost a more Bismarckian design shifts next period's benefits one for one from the flat to the earnings-related component, $\partial_{\theta}A = -\partial_{\theta}B = 1$, and the redistributive stake of the young dominates as in appendix \ref{app:esc:corner}. With it, a more Bismarckian design also raises the total the young will receive, since it dissipates less, which is what lets the design condition cross zero in the interior.

\subsection{The design choice under log preferences}\label{app:esc:log}
\begin{proposition}\label{prop:esc:separability}
Assume log-GHH preferences and no informal households, and let the design be chosen one period in advance under the cost \eqref{eq:esc:AB}. Then the political objective at $t$ is additively separable,
\begin{align}\label{eq:esc:separability}
    \mathcal{W}_t = \mathcal{A}_t\left(\tau_t,\theta_t\right) + \mathcal{B}_t\left(\theta_{t+1}\right) + \mathcal{C}_t\left(s_{t-1}\right),
\end{align}
where $\mathcal{C}_t$ depends on no policy. The tax chosen at $t$ depends on the design in force, $\theta_t$, and on neither $\theta_{t+1}$ nor $s_{t-1}$; the design chosen at $t$ maximizes $\mathcal{B}_t$ alone and depends on neither $\theta_t$ nor $s_{t-1}$.
\end{proposition}

\smalltitle{Proof of proposition \ref{prop:esc:separability}.} By \eqref{eq:esc:auxiliary:Thetah}, $\Theta_{h,t}$ is the product of a factor in $\tau_t$ and a factor in $(\tau_{t+1},\theta_{t+1})$, and $\tau_{t+1} = \mathcal{T}_{t+1}(\theta_{t+1})$ by the tax condition above, so by \eqref{eq:h} $\ln h_t$ is the sum of a term in $\tau_t$, a term in $\theta_{t+1}$ and $\frac{\alpha\xi}{1+\alpha\xi}\ln(s_{t-1}/\nu_t)$. The same holds for $\ln s_t = \frac{1+\xi}{\xi}\ln(h_t/\Gamma_h) + \ln\Gamma_{s,t}$ by \eqref{eq:s}, whose last term is a function of $\theta_{t+1}$ alone, and for $\ln\tilde{c}_{1,t}^i$, which by appendix \ref{app:EE:CRRA} is $\ln s_t$ plus the logarithm of a bracket in $\Gamma_{s,t}$, $\theta_{t+1}$ and $\tau_{t+1}$. The return $\ln R_{t+1}$ is affine in $\ln(s_t/\nu_{t+1})$ and $\ln h_{t+1}$, and $\ln h_{t+1}$ is $\ln\Theta_{h,t+1}$ plus a multiple of $\ln(s_t/\nu_{t+1})$, where $\Theta_{h,t+1}$ depends on $\tau_{t+1}$ and on the continuation $(\theta_{t+2},\tau_{t+2})$; by backward induction from the terminal period, in which no design is chosen, the design chosen at $t+1$ depends on no state, so the continuation is a function of $\theta_{t+1}$ alone and every term of the young's utility is a sum of terms in $\tau_t$, in $\theta_{t+1}$ and in $s_{t-1}$. Finally, $\ln c_{2,t}^i = (1-\alpha)\ln h_t + \ln\big[s_{t-1}^i/s_{t-1} + \frac{1-\alpha}{\alpha}\tau_t\left(A(\theta_t,\tau_t)\,y_i + B(\theta_t,\tau_t)\right)\big] + \alpha\ln(s_{t-1}/\nu_t) + \text{const}$, and $s_{t-1}^i/s_{t-1}$ is a function of $(\tau_t,\theta_t)$ by \eqref{eq:esc:auxiliary:si} at $t-1$. Every dependence of \eqref{eq:LOG:PolObj} on $\theta_{t+1}$ therefore arrives additively, apart from the terms in $(\tau_t,\theta_t)$ and in $s_{t-1}$, which is \eqref{eq:esc:separability}. The first order condition for the tax is $\partial\mathcal{A}_t/\partial\tau_t = 0$, and the design chosen at $t$ is the maximizer of $\mathcal{B}_t$.

The design chosen at $t$ is thus the same whatever system the electorate inherits and whatever the capital stock, which is what section \ref{sec:numerical} uses: each date's choice is a search over a grid of designs with no state, the tax at each candidate design is found as in proposition \ref{prop:LOG:PEE}, and the calibration condition \eqref{eq:esc:calibration} is a function of $\lambda$ alone. The property is specific to log preferences. Under CRRA preferences the powers do not separate, both $s_{t-1}$ and $\theta_t$ re-enter the design choice, and the equilibrium is a backward recursion over policy functions of the state $(s_{t-1},\theta_t)$ \parencite{BergGE26doc}.
